\documentclass[../../../main.tex]{subfiles}

\begin{document}
% Exercise 1.4:2(c)

Bergman suggested showing that $m_0 + 1$ is an upper bound of $\left\{\sum_{i=0}^{k} m_i10^{-i} \middle| k \geq 0 \right\}$ for this part, we might as well go for $\sum_{i=0}^{s} m_i 10^{-i} + 10^{-s}$ for any $s$ is an upper bound of this set.

\begin{proof}
	Fix any $s$, if $k \leq s$, it's easy to see that
	\[ 
		\sum_{i=0}^{s} m_i 10^{-i} + 10^{-s} \geq \sum_{i=0}^{k} m_i10^{-i},
	 \]
	since the additional terms between the $k$-th and $s$-th partial sums are nonnegative.

	If $k > s$, then
	\[
		\begin{aligned}
		\sum_{i=0}^{s} m_i 10^{-i} + 10^{-s}
		- \sum_{i=0}^{k} m_i 10^{-i}
		&= 10^{-s} - \sum_{i=s+1}^{k} m_i 10^{-i} \\
		&\geq 10^{-(s+1)} - \sum_{i=s+2}^{k} m_i 10^{-i} \\
		&\geq \cdots \\
		&\geq 10^{-k} \\
		&> 0.
		\end{aligned}
	\]
	In either case, we see that $\sum_{i=0}^{s} m_i 10^{-i} + 10^{-s}$ is an upper bound of $\left\{\sum_{i=0}^{k} m_i10^{-i} \middle| k \geq 0 \right\}$.


\end{proof}


\end{document}
